% ============================================================
% AUTOMATISCH GEGENEREERD BESTAND
% Niet handmatig aanpassen.
% ============================================================

\ifdefined\HCode
  % Geen afzonderlijke module-openingspagina in HTML.
\else
  \FVDmoduleopening
    {Functies in verandering}
    {In deze cursus bouwen we verder op functies die je al kent: veeltermfuncties, rationale en irrationale functies, exponentiële en logaritmische functies. We onderzoeken nu niet alleen hoe hun grafieken eruitzien, maar vooral **hoe ze veranderen**. Met afgeleiden bestuderen we stijgen en dalen, extrema, kromming en raaklijnen. Daarna breiden we onze gereedschapskist uit naar goniometrische en cyclometrische functies en leren we ook daar veranderingen nauwkeurig beschrijven. Maar analyse kijkt niet alleen naar verandering in één punt. We onderzoeken ook hoe vele kleine veranderingen samen een groter geheel vormen. Zo komen we stap voor stap bij oppervlakten en integralen terecht. Onderweg duiken dezelfde ideeën op in wetenschappelijke contexten zoals groei en verval, trillingen, optimalisatie en andere modellen. Zo wordt analyse een krachtig instrument om functies niet alleen te tekenen, maar echt te **begrijpen, onderzoeken en gebruiken**.}
    {%
    \FVDmodulethemetile{1}{Veranderingen in beeld}
    \FVDmodulethemetile{2}{Van oppervlakte naar integraal}
    \FVDmodulethemetile{3}{Functies met breuken en wortels}
    \FVDmodulethemetile{4}{Groei,verval en logaritmen}
    \FVDmodulethemetile{5}{Golven en periodiciteit}
    \FVDmodulethemetile{6}{Omgekeerde functies}
    \FVDmodulethemetile{7}{Integralen begrijpen}
    \FVDmodulethemetile{8}{Integreren met strategie}
    }
\fi
